\subsection{The Wonderful Exchange}\label{sec:exchange}

On 1 January, the eighth day of Christmas, the Church sings this antiphon at
Lauds:

\begin{studyframe}
\latin{O admirabile commercium: Creator generis humani, animatum corpus sumens,
de Virgine nasci dignatus est: et procedens homo sine semine, largitus est
nobis suam Deitatem.}
\medskip

How wonderful is the dealing here! The Maker of man hath taken to Himself a
living body, and hath been pleased to be born of the Virgin: He is come forth
conceived without seed, and hath made us partakers of His Divine nature.\par
\end{studyframe}

\latin{Commercium} is a word of the market: trade, a dealing, an exchange. Two
parties meet, and each gives the other what he has. In this exchange the Maker
of man takes a living body, and gives his Godhead. St Leo the Great, preaching
on the Passion, says that Christ had come into this world ``a rich and merciful
Merchant from the skies, and by a wondrous exchange had entered into a bargain
of salvation with us, receiving ours and giving His, honour for insults,
salvation for pain, life for death''.

The antiphon's last words echo St Peter, who writes that God has given us
``most great and precious promises: that by these you may be made partakers of
the divine nature''. St Paul tells what it cost the one who gave:

\begin{quote}
For let this mind be in you, which was also in Christ Jesus: Who being in the
form of God, thought it not robbery to be equal with God: But emptied himself,
taking the form of a servant, being made in the likeness of men, and in habit
found as a man. He humbled himself, becoming obedient unto death, even to the
death of the cross.\par
{\smallskip\noindent\small ---Philippians 2:5--8\par}
\end{quote}

\subsubsection*{O new commingling}

St Gregory Nazianzen, preaching on the Birthday of Christ, piles one wonder on
another:

\begin{quote}
O new commingling; O strange conjunction; the Self-Existent comes into being,
the Uncreate is created, That which cannot be contained is contained, by the
intervention of an intellectual soul, mediating between the Deity and the
corporeity of the flesh. And He Who gives riches becomes poor, for He assumes
the poverty of my flesh, that I may assume the richness of His Godhead. He that
is full empties Himself, for He empties Himself of His glory for a short while,
that I may have a share in His Fulness.\par
{\smallskip\noindent\small ---St Gregory Nazianzen, Oration 38, on the Birthday of Christ\par}
\end{quote}

St Athanasius of Alexandria wrote:

\begin{quote}
For He was made man that we might be made God; and He manifested Himself by a
body that we might receive the idea of the unseen Father; and He endured the
insolence of men that we might inherit immortality.\par
{\smallskip\noindent\small ---St Athanasius, On the Incarnation of the Word\par}
\end{quote}

To be ``made God'' is St Peter's ``partakers of the divine nature'' said as sharply
as it can be said: we are given, by grace, a share in the life of the One who
took our nature, so that what is his becomes ours as truly as what was ours
became his.

\subsubsection*{The Bread hungers}

St Augustine, in a sermon for Christmas Day, lets the paradoxes run until they
reach the Cross. The Maker of man was made man, \latin{Homo factus, hominis
factor}, that he who rules the stars might suck at the breast; and then,
\latin{ut esuriret panis, ut sitiret fons, dormiret lux, ab itinere via
fatigaretur}: that the Bread might hunger, the Fountain thirst, the Light sleep,
the Way be wearied by the journey. The list does not stop at the manger. The
Truth is accused by false witnesses, the Judge of the living and the dead is
judged by a mortal judge, the Cluster of grapes is crowned with thorns, the
Foundation is hung upon the wood; and it ends, \latin{vita moreretur}, that Life
might die.

The Bread that hungers is the Child born in the House of Bread. The Way that is
wearied is the Child born by the road. Augustine sees the whole exchange in the
Nativity, down to the thorns and the wood.

In another Christmas sermon he quotes a psalm, ``Truth is sprung out of the
earth: and justice hath looked down from heaven'', and says of it,
\latin{Veritas cui cœlum non sufficit, de terra orta est, ut in præsepio
poneretur}: ``Truth, which heaven is not wide enough to hold, sprang out of the
earth, to be laid in a manger.'' Then he turns on his hearers:
\latin{Expergiscere, homo: pro te Deus factus est homo}, ``Awake, O man: for you
God was made man.'' Advent began with St Paul's ``it is now the hour for us to
rise from sleep''; at Christmas Augustine wakes his hearers with St Paul again,
``Rise, thou that sleepest, and arise from the dead: and Christ shall enlighten
thee.'' In one short sentence after another he tells them what they would
have been without him: dead for ever, had he not been born in time; never set
free from the flesh of sin, had he not taken its likeness. The last is the
shortest, and it ends on the verb the Church has said to Christ all through
Advent: \latin{Perisses, nisi venisset}, you would have perished, had he not
come.

On Christmas Day, and again on 1 January, say the antiphon \latin{O admirabile
commercium} (Appendix~\ref{sec:prayers}), slowly, and stop at its last words.
At Christmas Matins the Church reads St Leo's first sermon on the Nativity, and
the last of its lessons is a charge to every baptized man and woman. Read it aloud at home on
Christmas Day:

\begin{quote}
Learn, O Christian, how great thou art, who hast been made partaker of the
Divine nature, \dots\ and fall not again by corrupt conversation into the
beggarly elements above which thou art lifted. Remember Whose Body it is
Whereof thou art made a member, and Who is its Head\dots\par
{\smallskip\noindent\small ---St Leo the Great, Sermon 21 on the Nativity\par}
\end{quote}
